\documentclass[a4paper,12pt]{article}

\usepackage[T2A]{fontenc}
\usepackage[utf8]{inputenc}
\usepackage[russian]{babel}

\usepackage{amsmath,amssymb}
\usepackage{icomma}

\usepackage{paratype}
\renewcommand{\familydefault}{\sfdefault}
\usepackage{microtype}
\usepackage{setspace}
\onehalfspacing
\usepackage{parskip}
\usepackage{enumitem}

\usepackage{geometry}
\geometry{
  a4paper,
  left=20mm,
  right=20mm,
  top=20mm,
  bottom=22mm
}

\usepackage{xcolor}
\usepackage{tikz}
\usetikzlibrary{calc,arrows.meta,positioning,shapes.geometric}
\usepackage{tkz-euclide}

\usepackage{tabularx,booktabs,array}
\usepackage{fancyhdr}
\usepackage{titlesec}

\definecolor{accent}{HTML}{008080}
\definecolor{darkaccent}{HTML}{004D4D}
\definecolor{textgray}{HTML}{2F3333}
\definecolor{softgray}{HTML}{6F7777}
\definecolor{lightaccent}{HTML}{D9EEEE}

\color{textgray}

\pagestyle{fancy}
\fancyhf{}
\fancyhead[L]{\small\color{softgray}Тригонометрические уравнения}
\fancyhead[R]{\small\color{softgray}Ярослав Гаврилов}
\fancyfoot[C]{\small\color{softgray}\thepage}
\setlength{\headheight}{25pt}
\renewcommand{\headrulewidth}{0.3pt}
\renewcommand{\headrule}{%
  \hbox to\headwidth{%
    \color{lightaccent}\leaders\hrule height \headrulewidth\hfill
  }%
}

\titleformat{\section}
  {\Large\bfseries\color{darkaccent}}
  {}
  {0pt}
  {}

\titleformat{\subsection}
  {\large\bfseries\color{accent}}
  {}
  {0pt}
  {}

\titlespacing*{\section}{0pt}{1.4em}{0.65em}
\titlespacing*{\subsection}{0pt}{1.0em}{0.4em}

\tikzset{
  flowStart/.style={
    ellipse,
    draw=darkaccent,
    line width=0.9pt,
    align=center,
    minimum width=42mm,
    minimum height=11mm,
    inner sep=4pt,
    font=\small
  },
  flowAction/.style={
    rounded corners=3pt,
    rectangle,
    draw=accent,
    line width=0.9pt,
    align=center,
    text width=49mm,
    minimum height=13mm,
    inner sep=5pt,
    font=\small
  },
  flowDecision/.style={
    diamond,
    aspect=2.2,
    draw=darkaccent,
    line width=0.9pt,
    align=center,
    text width=34mm,
    inner sep=2pt,
    font=\small
  },
  flowArrow/.style={
    -{Stealth[length=3mm,width=2mm]},
    line width=0.9pt,
    draw=darkaccent
  },
  flowLabel/.style={
    font=\small,
    fill=white,
    inner sep=1.5pt
  }
}

\newcommand{\Z}{\mathbb{Z}}

\setlist[itemize]{
  leftmargin=1.5em,
  itemsep=0.22em,
  topsep=0.3em
}

\setlist[enumerate]{
  leftmargin=1.7em,
  itemsep=0.4em,
  topsep=0.35em
}

\renewcommand{\arraystretch}{1.3}

\begin{document}

\begin{titlepage}
\thispagestyle{empty}

\begin{center}
\vspace*{23mm}

{\color{accent}\rule{0.22\linewidth}{1.2pt}}

\vspace{11mm}

{\Huge\bfseries\color{darkaccent}Чек-лист}

\vspace{5mm}

{\LARGE\bfseries Тригонометрические уравнения}

\vspace{4mm}

{\large
Единичная окружность, табличные значения,\\
арксинус и арккосинус, частные и общие случаи
}

\vspace{20mm}

{\large Ярослав Гаврилов}

\vspace{4mm}

{\normalsize Подготовка к ЕГЭ}

\vfill

{\large
Лёвин Артём Александрович\\
эксклюзивно для Ярослава Гаврилова
}

\vspace{4mm}

{\normalsize \today}

\vspace{26mm}

\begin{tikzpicture}[remember picture,overlay]
  \draw[
    accent,
    line width=1.1pt,
    opacity=0.65
  ]
  ([xshift=18mm,yshift=18mm]current page.south west)
  .. controls
  ([xshift=70mm,yshift=32mm]current page.south west)
  and
  ([xshift=-80mm,yshift=8mm]current page.south east)
  ..
  ([xshift=-18mm,yshift=23mm]current page.south east);
\end{tikzpicture}

\end{center}
\end{titlepage}

\section{1. Единичная окружность: что означают синус и косинус}

Тригонометрическая окружность имеет центр в начале координат и радиус, равный~1.
Если точка \(M\) соответствует углу \(\alpha\), то

\[
\boxed{M(\cos\alpha;\,\sin\alpha)}.
\]

Следовательно,

\[
\boxed{x_M=\cos\alpha},
\qquad
\boxed{y_M=\sin\alpha}.
\]

Положительные углы отсчитываются против часовой стрелки, отрицательные --- по часовой.

\begin{center}
\begin{tikzpicture}[scale=2.2,>=Stealth]
  \tkzDefPoint(0,0){O}
  \tkzDefPoint(1,0){A}
  \tkzDefPoint(0.766,0.643){M}

  \draw[->,line width=0.7pt] (-1.35,0)--(1.45,0) node[right] {\(x\)};
  \draw[->,line width=0.7pt] (0,-1.35)--(0,1.45) node[above] {\(y\)};

  \tkzDrawCircle[color=accent,line width=1.1pt](O,A)
  \tkzDrawSegment[color=darkaccent,line width=1pt](O,M)

  \draw[dashed,softgray] (M)--(0.766,0);
  \draw[dashed,softgray] (M)--(0,0.643);
  \draw[->,accent,line width=0.9pt]
    (0.38,0) arc[start angle=0,end angle=40,radius=0.38];

  \tkzDrawPoints[color=darkaccent,fill=darkaccent,size=3](M)
  \fill (1,0) circle (0.018);
  \fill (0,1) circle (0.018);
  \fill (-1,0) circle (0.018);
  \fill (0,-1) circle (0.018);

  \node[below right] at (1,0) {\((1;0)\)};
  \node[above right] at (0,1) {\((0;1)\)};
  \node[below left] at (-1,0) {\((-1;0)\)};
  \node[below right] at (0,-1) {\((0;-1)\)};
  \node[above right] at (M) {\(M(\cos\alpha;\sin\alpha)\)};
  \node at (0.48,0.16) {\(\alpha\)};
\end{tikzpicture}
\end{center}

\subsection{Ключевые наблюдения}

\begin{itemize}
  \item значение \(\sin\alpha\) определяется вертикальной координатой точки;
  \item значение \(\cos\alpha\) определяется горизонтальной координатой точки;
  \item после полного оборота точка возвращается на прежнее место;
  \item полный оборот равен
  \[
  2\pi=360^\circ.
  \]
\end{itemize}

Поэтому углы

\[
\alpha=\alpha_0+2\pi n,
\qquad n\in\Z,
\]

задают одну и ту же точку окружности при всех целых \(n\).

Кроме того,

\[
-1\le \sin\alpha\le 1,
\qquad
-1\le \cos\alpha\le 1.
\]

Это сразу даёт важное ограничение: уравнения \(\sin x=a\) и \(\cos x=a\) имеют действительные решения только при \(|a|\le1\).

\section{2. Табличные значения}

Основные рабочие углы удобно знать и в градусах, и в радианах.

\begin{center}
\begin{tabularx}{0.98\linewidth}{
  >{\centering\arraybackslash}X
  >{\centering\arraybackslash}X
  >{\centering\arraybackslash}X
  >{\centering\arraybackslash}X
  >{\centering\arraybackslash}X
  >{\centering\arraybackslash}X
}
\toprule
\textbf{Угол} & \(0^\circ\) & \(30^\circ\) & \(45^\circ\) & \(60^\circ\) & \(90^\circ\) \\
\midrule
\textbf{Радианы} & \(0\) & \(\dfrac{\pi}{6}\) & \(\dfrac{\pi}{4}\) & \(\dfrac{\pi}{3}\) & \(\dfrac{\pi}{2}\) \\[1mm]
\(\sin\alpha\) & \(0\) & \(\dfrac12\) & \(\dfrac{\sqrt2}{2}\) & \(\dfrac{\sqrt3}{2}\) & \(1\) \\[2mm]
\(\cos\alpha\) & \(1\) & \(\dfrac{\sqrt3}{2}\) & \(\dfrac{\sqrt2}{2}\) & \(\dfrac12\) & \(0\) \\
\bottomrule
\end{tabularx}
\end{center}

Таблицу можно быстро восстановить по шаблону

\[
\sin:\quad
\frac{\sqrt0}{2},\
\frac{\sqrt1}{2},\
\frac{\sqrt2}{2},\
\frac{\sqrt3}{2},\
\frac{\sqrt4}{2},
\]

\[
\cos:\quad
\frac{\sqrt4}{2},\
\frac{\sqrt3}{2},\
\frac{\sqrt2}{2},\
\frac{\sqrt1}{2},\
\frac{\sqrt0}{2}.
\]

\section{3. Арксинус и арккосинус}

Для \(a\in[-1;1]\) число \(\arcsin a\) --- это единственный угол из промежутка

\[
\left[-\frac{\pi}{2};\frac{\pi}{2}\right],
\]

синус которого равен \(a\):

\[
\sin(\arcsin a)=a.
\]

Аналогично, \(\arccos a\) --- единственный угол из промежутка

\[
[0;\pi],
\]

косинус которого равен \(a\):

\[
\cos(\arccos a)=a.
\]

Важно: \(\arcsin a\) и \(\arccos a\) задают \emph{главные значения}. Они сами по себе ещё не описывают все решения тригонометрического уравнения.

\subsection{Примеры}

\[
\arcsin\frac12=\frac{\pi}{6},
\qquad
\arcsin\left(-\frac12\right)=-\frac{\pi}{6},
\]

\[
\arccos\frac12=\frac{\pi}{3},
\qquad
\arccos\left(-\frac12\right)=\frac{2\pi}{3}.
\]

Если значение нетабличное, например \(\sin\alpha=0{,}8\), главный угол удобно обозначить как \(\arcsin 0{,}8\).

\section{4. Частные случаи}

Частные случаи соответствуют значениям \(-1\), \(0\), \(1\). Их удобно восстанавливать непосредственно по единичной окружности.

\subsection{Синус}

\[
\boxed{\sin x=0\quad\Longrightarrow\quad x=\pi n,
\qquad n\in\Z.}
\]

\[
\boxed{\sin x=1\quad\Longrightarrow\quad x=\frac{\pi}{2}+2\pi n,
\qquad n\in\Z.}
\]

\[
\boxed{\sin x=-1\quad\Longrightarrow\quad x=-\frac{\pi}{2}+2\pi n,
\qquad n\in\Z.}
\]

Для последнего случая эквивалентна запись

\[
x=\frac{3\pi}{2}+2\pi n,
\qquad n\in\Z.
\]

\subsection{Косинус}

\[
\boxed{\cos x=0\quad\Longrightarrow\quad x=\frac{\pi}{2}+\pi n,
\qquad n\in\Z.}
\]

\[
\boxed{\cos x=1\quad\Longrightarrow\quad x=2\pi n,
\qquad n\in\Z.}
\]

\[
\boxed{\cos x=-1\quad\Longrightarrow\quad x=\pi+2\pi n,
\qquad n\in\Z.}
\]

\section{5. Общий случай: \(\sin x=a\)}

Если \(|a|>1\), действительных решений нет.

Если \(|a|\le1\), полное множество решений удобно записывать двумя сериями:

\[
\boxed{
\begin{aligned}
x&=\arcsin a+2\pi n,\\
x&=\pi-\arcsin a+2\pi n,
\end{aligned}
\qquad n\in\Z.
}
\]

При \(a=\pm1\) эти две серии описывают одно и то же множество решений, поэтому в частных случаях запись можно сократить.

\subsection{Пример 1: \(\sin x=\frac12\)}

\[
\arcsin\frac12=\frac{\pi}{6}.
\]

Поэтому

\[
x=\frac{\pi}{6}+2\pi n
\]

или

\[
x=\pi-\frac{\pi}{6}+2\pi n
=\frac{5\pi}{6}+2\pi n.
\]

Итак,

\[
\boxed{
x=\frac{\pi}{6}+2\pi n
\quad\text{или}\quad
x=\frac{5\pi}{6}+2\pi n,
\qquad n\in\Z.
}
\]

\subsection{Пример 2: \(\sin x=-\frac12\)}

\[
\arcsin\left(-\frac12\right)=-\frac{\pi}{6}.
\]

Следовательно,

\[
x=-\frac{\pi}{6}+2\pi n
\]

или

\[
x=\pi+\frac{\pi}{6}+2\pi n
=\frac{7\pi}{6}+2\pi n.
\]

Эквивалентная компактная формула для общего случая:

\[
\boxed{x=(-1)^n\arcsin a+\pi n,
\qquad n\in\Z.}
\]

Для отбора корней по окружности часто удобнее сохранять две отдельные серии: так видны обе точки и оба направления рассуждения.

\section{6. Общий случай: \(\cos x=a\)}

Если \(|a|>1\), действительных решений нет.

Если \(|a|\le1\), точки с одинаковой горизонтальной координатой симметричны относительно оси \(Ox\). Поэтому

\[
\boxed{x=\pm\arccos a+2\pi n,
\qquad n\in\Z.}
\]

\subsection{Пример 1: \(\cos x=\frac12\)}

\[
\arccos\frac12=\frac{\pi}{3},
\]

поэтому

\[
\boxed{x=\pm\frac{\pi}{3}+2\pi n,
\qquad n\in\Z.}
\]

\subsection{Пример 2: \(\cos x=-\frac12\)}

\[
\arccos\left(-\frac12\right)=\frac{2\pi}{3}.
\]

Следовательно,

\[
\boxed{x=\pm\frac{2\pi}{3}+2\pi n,
\qquad n\in\Z.}
\]

В пределах одного оборота это соответствует углам

\[
\frac{2\pi}{3}
\qquad\text{и}\qquad
\frac{4\pi}{3}.
\]

\section{7. Если аргумент функции сложнее, чем \(x\)}

Полезный приём: сначала рассматривайте весь аргумент функции как один угол, решайте простейшее тригонометрическое уравнение относительно этого угла, затем выражайте \(x\).

Например,

\[
\sin3x=1.
\]

Сначала решаем относительно аргумента \(3x\):

\[
3x=\frac{\pi}{2}+2\pi n,
\qquad n\in\Z.
\]

Делим всю правую часть на~3:

\[
\boxed{x=\frac{\pi}{6}+\frac{2\pi n}{3},
\qquad n\in\Z.}
\]

Типичная ошибка --- разделить на коэффициент только первый член справа. Проверяйте деление всей суммы:

\[
\frac{\frac{\pi}{2}+2\pi n}{3}
=
\frac{\pi}{6}+\frac{2\pi n}{3}.
\]

\clearpage
\thispagestyle{empty}

\begin{center}
{\LARGE\bfseries\color{darkaccent}
Алгоритм: простейшее тригонометрическое уравнение
}

\vspace{9mm}

\begin{tikzpicture}[node distance=8mm and 15mm]

\node[flowStart] (start)
{Получено уравнение\\
\(\sin u=a\) или \(\cos u=a\)};

\node[flowDecision, below=of start] (range)
{\(|a|\le1\)?};

\node[flowStart, right=15mm of range, minimum width=45mm] (none)
{Действительных\\решений нет};

\node[flowDecision, below=13mm of range] (special)
{\(a\in\{-1,0,1\}\)?};

\node[flowAction, left=15mm of special, text width=43mm] (circle)
{Найти подходящие точки по единичной окружности и записать периодические решения};

\node[flowDecision, below=13mm of special] (func)
{Какая функция?};

\node[flowAction, below left=14mm and 8mm of func, text width=45mm] (sinbranch)
{Для \(\sin u=a\):\\[1mm]
\(u=\arcsin a+2\pi n\)\\
или\\
\(u=\pi-\arcsin a+2\pi n\)};

\node[flowAction, below right=14mm and 8mm of func, text width=45mm] (cosbranch)
{Для \(\cos u=a\):\\[1mm]
\(u=\pm\arccos a+2\pi n\)};

\node[flowAction, below=15mm of $(sinbranch)!0.5!(cosbranch)$] (solve)
{Если \(u\neq x\), решить полученные\\
уравнения относительно \(x\)};

\node[flowAction, below=of solve] (check)
{Указать \(n\in\Z\) и проверить ограничения};

\node[flowStart, below=of check] (finish)
{Записать ответ};

\draw[flowArrow] (start)--(range);
\draw[flowArrow] (range.east)--node[flowLabel,above]{нет}(none.west);
\draw[flowArrow] (range.south)--node[flowLabel,right]{да}(special.north);
\draw[flowArrow] (special.west)--node[flowLabel,above]{да}(circle.east);
\draw[flowArrow] (special.south)--node[flowLabel,right]{нет}(func.north);

\draw[flowArrow]
  (func.south west)--node[flowLabel,above left]{синус}(sinbranch.north);
\draw[flowArrow]
  (func.south east)--node[flowLabel,above right]{косинус}(cosbranch.north);

\draw[flowArrow] (sinbranch.south) |- (solve.west);
\draw[flowArrow] (cosbranch.south) |- (solve.east);
\draw[flowArrow] (circle.south) |- ([xshift=-24mm]solve.west) -- (solve.west);
\draw[flowArrow] (solve)--(check);
\draw[flowArrow] (check)--(finish);

\end{tikzpicture}
\end{center}

\clearpage

\section{8. Как формулы превращают сложное уравнение в простейшее}

В более содержательных задачах сначала нужно увидеть знакомую тригонометрическую конструкцию, применить формулу и получить простейшее уравнение.

Одна из основных формул:

\[
\boxed{\sin2\alpha=2\sin\alpha\cos\alpha.}
\]

Рассмотрим

\[
8\sin2x\cos2x=4.
\]

Выделяем удвоенное произведение:

\[
4\bigl(2\sin2x\cos2x\bigr)=4.
\]

По формуле двойного угла

\[
2\sin2x\cos2x=\sin4x,
\]

поэтому

\[
4\sin4x=4,
\]

\[
\sin4x=1.
\]

Теперь используем частный случай:

\[
4x=\frac{\pi}{2}+2\pi n,
\qquad n\in\Z.
\]

Отсюда

\[
\boxed{x=\frac{\pi}{8}+\frac{\pi n}{2},
\qquad n\in\Z.}
\]

Рабочая цепочка:

\[
\text{узнать формулу}
\;\longrightarrow\;
\text{преобразовать выражение}
\]

\[
\text{получить простейшее уравнение}
\;\longrightarrow\;
\text{решить его}.
\]

\section{9. Типичные ошибки и самопроверка}

\begin{enumerate}
  \item Помните координатный смысл:
  \[
  x_M=\cos\alpha,
  \qquad
  y_M=\sin\alpha.
  \]

  \item Всегда проверяйте условие
  \[
  |a|\le1
  \]
  в уравнениях \(\sin x=a\) и \(\cos x=a\).

  \item Не теряйте периодичность. Одна точка на окружности соответствует бесконечному множеству углов.

  \item Для \(\sin x=0\) и \(\cos x=0\) удобно использовать сокращённые записи с периодом \(\pi\):
  \[
  \sin x=0\Rightarrow x=\pi n,
  \qquad
  \cos x=0\Rightarrow x=\frac{\pi}{2}+\pi n.
  \]

  \item Если аргумент имеет вид \(kx\), после решения уравнения относительно \(kx\) делите на \(k\) всю правую часть.

  \item При использовании \(\arcsin\) и \(\arccos\) различайте главное значение и полное множество решений.

  \item Для самопроверки подставьте один-два простых значения \(n\), например \(0\) и \(1\), и убедитесь по окружности, что полученные углы действительно дают нужный синус или косинус.
\end{enumerate}

\section{10. Короткая памятка}

\[
\boxed{M(\cos\alpha;\sin\alpha)}
\]

\[
\boxed{
\sin x=a\quad\Longrightarrow\quad
\begin{cases}
x=\arcsin a+2\pi n,\\
x=\pi-\arcsin a+2\pi n,
\end{cases}
\quad |a|\le1,
\quad n\in\Z.
}
\]

\[
\boxed{
\cos x=a\quad\Longrightarrow\quad
x=\pm\arccos a+2\pi n,
\quad |a|\le1,
\quad n\in\Z.
}
\]

\clearpage

\section{Тренировочный блок --- Путь А}

Решите уравнения. Решения в тренировочном блоке не записывайте в сокращённом виде до тех пор, пока не проверите период и все серии корней.

\begin{enumerate}
  \item
  \[
  \sin x=0.
  \]

  \vspace{17mm}

  \item
  \[
  \cos x=-\frac12.
  \]

  \vspace{20mm}

  \item
  \[
  \sin3x=1.
  \]

  \vspace{20mm}

  \item
  \[
  \cos2x=\frac{\sqrt2}{2}.
  \]

  \vspace{22mm}

  \item
  \[
  8\sin2x\cos2x=4.
  \]

  \vspace{24mm}
\end{enumerate}

\begin{center}
\color{softgray}
Сначала определите тип уравнения. Если требуется, преобразуйте выражение до простейшей тригонометрической формы.
\end{center}

\clearpage

\section{Ответы}

\begin{center}
\begin{tabularx}{0.96\linewidth}{
  >{\centering\arraybackslash}p{12mm}
  X
}
\toprule
\textbf{№} & \textbf{Ответ} \\
\midrule
1 & \(\displaystyle x=\pi n,\qquad n\in\Z.\) \\
\addlinespace[3mm]
2 & \(\displaystyle x=\pm\frac{2\pi}{3}+2\pi n,\qquad n\in\Z.\) \\
\addlinespace[3mm]
3 & \(\displaystyle x=\frac{\pi}{6}+\frac{2\pi n}{3},\qquad n\in\Z.\) \\
\addlinespace[3mm]
4 & \(\displaystyle x=\pm\frac{\pi}{8}+\pi n,\qquad n\in\Z.\) \\
\addlinespace[3mm]
5 & \(\displaystyle x=\frac{\pi}{8}+\frac{\pi n}{2},\qquad n\in\Z.\) \\
\bottomrule
\end{tabularx}
\end{center}

\vfill

\begin{center}
\begin{tikzpicture}
  \draw[
    accent,
    line width=1pt,
    opacity=0.6
  ]
  (0,0)
  .. controls (3,0.75) and (8,-0.55) ..
  (12,0.15);
\end{tikzpicture}

\vspace{2mm}

{\small\color{softgray}Лёвин Артём Александрович}
\end{center}

\end{document}